% 結論
本研究では,強相関電子系におけるモット転移の性質を明らかにするため,ghost Gutzwiller近似を用いた数値計算を実装し,以下の知見を得た.

\begin{itemize}
    \item \textbf{モット転移とヒステリシスの観測}:
    単一バンドハバードモデルにおいて,相互作用パラメータの増大に伴い準粒子重みがゼロとなる金属絶縁体転移を明確に観測した.

    また,昇順・降順スイープを行うことで,一次相転移に特有のヒステリシスループを捉えることに成功した.

    \item \textbf{厳密解との定量的な一致}:
    無限次元Bethe格子における転移点の臨界値は,動的平均場理論による厳密な計算結果と極めて良い一致を示した.

    これにより,少数のゴースト軌道を用いて変分空間を拡張する本手法が,モット転移の物理を高精度に記述できることが証明された.

    \item \textbf{計算コストの劇的な削減と手法の優位性}:
    動的平均場理論に比べて,本手法は無限個の電子の浴の役割をたった2つの空間ゴースト軌道に押し込めることで,不純物問題を劇的に単純化している.

    厳密対角化を用いて一瞬で解くことができるため,極めて低い計算コストで高精度な結果を得られることが確認された.

    \item \textbf{今後の展望}:
    本手法の計算コストの低さと精度の高さを活かし,今後は多軌道系やより複雑な格子構造を持つ現実の物質の解析へ応用することが期待される.
\end{itemize}